% !TeX program = xelatex
% !TeX encoding = UTF-8 Unicode
% !TeX TXS-program:compile = txs:///xelatex
% !TeX TXS-program:quick = txs:///xelatex | txs:///view
\documentclass[11pt,a4paper]{ctexart}
\usepackage[margin=2.1cm]{geometry}
\usepackage{amsmath}
\usepackage{booktabs}
\usepackage{tabularx}
\usepackage{array}
\usepackage{enumitem}
\usepackage{hyperref}
\usepackage{xcolor}
\usepackage{listings}

\hypersetup{hidelinks}
\setlength{\parindent}{2em}
\setlength{\parskip}{0.25em}
\renewcommand{\arraystretch}{1.12}
\setlist[itemize]{leftmargin=2em,itemsep=0.1em,topsep=0.2em}
\setlist[enumerate]{leftmargin=2.2em,itemsep=0.1em,topsep=0.2em}
\lstset{
  basicstyle=\ttfamily\small,
  breaklines=true,
  columns=fullflexible,
  frame=single,
  backgroundcolor=\color{black!3}
}

\title{\textbf{面向 Ansys Maxwell 的工业软件智能体原型}\\
\large 自然语言驱动的仿真规格生成、软件执行与反馈修正}
\author{阶段汇报精简版}
\date{2026年4月29日}

\begin{document}
\maketitle

\begin{abstract}
本项目以 Ansys Maxwell 为对象，构建了一个工业软件智能体原型。用户输入自然语言需求后，系统将需求整理为仿真规格和执行计划，生成 PyAEDT 脚本，在本地调用 Maxwell 完成建模、求解和工程文件保存，并对仿真结果进行约束判定。当关键约束不满足时，系统会把结果和失败原因反馈给大模型，由大模型生成下一轮修改方案，再重新仿真验证。当前系统已验证二维电磁铁、二维平行板电容器、二维变压器概念模型和二维电感器概念模型四类任务，能够生成真实的 \texttt{.aedt} 工程文件和结构化结果。
\end{abstract}

\noindent\textbf{关键词：} 工业软件智能体；Ansys Maxwell；PyAEDT；大模型；仿真自动化

\section{问题与目标}

工业仿真软件的难点不只是操作复杂，更在于用户需要把工程语言转换成软件可执行的建模、激励、边界和求解流程。例如“做一个24V直流电磁铁，气隙2mm，电流不超过2A”这句话，对软件来说并不是一个可直接执行的任务。它需要先被转换为几何参数、线圈参数、材料、激励、求解器设置和结果判定规则。

本项目的目标是打通一条可执行链路：

\begin{center}
自然语言需求 $\rightarrow$ 仿真规格 $\rightarrow$ PyAEDT 脚本 $\rightarrow$ Maxwell 求解 $\rightarrow$ \texttt{.aedt} 工程与结果判定
\end{center}

这里的大模型不是直接生成 \texttt{.aedt} 文件，而是负责需求理解、仿真规格生成和失败后的设计修正；本地执行器负责脚本校验、Maxwell 调用、结果提取和约束判断。这样既利用了大模型的理解能力，也保留了工业软件执行的确定性。

\section{方法}

\subsection{统一中间表示}

系统先把用户需求整理成统一的仿真规格。仿真规格可以表示为：

\[
\Sigma = \{G, M, E, B, S, Q, C\}
\]

其中 $G$ 表示几何，$M$ 表示材料，$E$ 表示激励，$B$ 表示边界条件，$S$ 表示求解器，$Q$ 表示需要提取的结果，$C$ 表示用户提出的约束。随后系统生成执行计划：

\[
P = \{a_1, a_2, \ldots, a_n\}
\]

每个 $a_i$ 对应一个 Maxwell 操作，例如创建几何、赋材料、施加电流或电压、建立求解设置、运行求解和导出结果。最终执行计划被转换为 PyAEDT 脚本，由本地 Maxwell 执行并保存 \texttt{.aedt} 工程。

\subsection{脚本执行与安全控制}

自动生成脚本存在不稳定风险。系统会先做静态检查，确认脚本具备入口函数、Maxwell 调用和工程保存逻辑，同时限制危险操作。脚本不会在主进程中直接执行，而是在独立 Python 进程中运行。这样即使脚本或 Maxwell 退出异常，也不会拖死主程序。

Student 版 Maxwell 对并发较敏感，因此网页端采用单任务运行策略，避免多个仿真任务同时启动造成工程锁冲突或进程残留。

\subsection{结果判定与反馈修正}

系统不把“Maxwell 跑完”直接等同于“需求满足”。每次仿真后，系统会读取输出并逐条判断约束。以电磁铁为例，若用户要求 24V 且电流不超过 2A，则系统会估算线圈电阻 $R_{\mathrm{coil}}$，并计算：

\[
I_{\mathrm{est}} = \frac{V}{R_{\mathrm{coil}}}
\]

只有当 $I_{\mathrm{est}} \leq I_{\max}$ 时，才认为供电电压和电流约束同时满足。若不满足，系统会把当前结果、失败原因和用户原始需求反馈给大模型，让大模型生成下一轮修改方案，再重新仿真。

\section{实现效果}

当前系统已验证四类 Maxwell 任务族。表 \ref{tab:results} 给出最新实验结果。

\begin{table}[h]
\centering
\caption{已验证任务与实验结果}
\label{tab:results}
\begin{tabularx}{\textwidth}{>{\raggedright\arraybackslash}p{3.0cm}X>{\centering\arraybackslash}p{1.5cm}}
\toprule
任务 & 关键结果 & 评价 \\
\midrule
二维电磁铁 & 最大磁密 $9.866\times10^{-5}$ T；线圈电阻 13.41 $\Omega$；24V 下估算电流 1.79 A，满足不超过 2A 的约束 & 通过 \\
二维平行板电容器 & 电容 198.18 pF；最大电场约 $3.70\times10^{5}$ V/m & 通过 \\
二维变压器概念模型 & 原边 10000 V，副边目标 220 V；匝比 0.022；估算副边电压 220 V & 通过 \\
二维电感器概念模型 & 估算电感 0.326 H；最大磁密 $9.029\times10^{-5}$ T & 通过 \\
\bottomrule
\end{tabularx}
\end{table}

其中，电磁铁任务已经验证了反馈修正能力。首轮设计若导致 24V 下电流超过上限，系统会把失败原因反馈给大模型，第二轮调整线圈设计后重新运行，最终得到满足电压和电流联合约束的方案。

\section{复现方式}

复现实验需要本地安装 Ansys Maxwell / AEDT，并配置 Python 3.12 环境和云端模型 API。项目根目录为：

\begin{lstlisting}[language={}]
F:\maxwell_agent_project
\end{lstlisting}

首先检查 Maxwell 和大模型接口：

\begin{lstlisting}[language=bash]
.\.venv\Scripts\python.exe -m maxwell_agent.cli probe-env
.\.venv\Scripts\python.exe -m maxwell_agent.cli smoke-llm
\end{lstlisting}

运行电磁铁实验：

\begin{lstlisting}[language=bash]
.\.venv\Scripts\python.exe -m maxwell_agent.cli run "做一个24V直流电磁铁，气隙2mm，电流不超过2A，尽量提高吸力，外形不要太大。"
\end{lstlisting}

运行电容器实验：

\begin{lstlisting}[language=bash]
.\.venv\Scripts\python.exe -m maxwell_agent.cli run "做一个二维平行板电容器，板间距1mm，板宽20mm，施加100V电压，求电场强度和电容。"
\end{lstlisting}

每次运行都会在 \texttt{workspace} 下生成独立目录。判断是否成功主要看三个文件：

\begin{itemize}
  \item \texttt{*.aedt}：Maxwell 工程文件是否生成。
  \item \texttt{outputs.json}：仿真结果是否导出。
  \item \texttt{evaluation.json}：总体评价是否为 \texttt{passed}。
\end{itemize}

\section{局限与下一步}

当前系统已经不是单一电磁铁脚本，但仍然是原型阶段。稳定验证的主要是四类二维或概念模型，尚不能覆盖任意 Maxwell 问题。变压器和电感器目前偏概念验证，还没有完整覆盖损耗、温升、复杂绕组、三维结构和瞬态耦合等工程细节。

下一步需要继续扩展 Maxwell 操作原语、后处理能力和反馈修正策略。重点不是继续增加演示页面，而是让更多任务能够从自然语言稳定落到真实 Maxwell 工程，并让系统对“哪些条件满足、哪些条件不满足”给出更可靠的工程判断。

\section{结论}

本项目完成了一条从自然语言到 Maxwell 仿真执行结果的智能体链路。系统能够生成真实 \texttt{.aedt} 工程文件，提取仿真结果，并对用户的数值约束进行判断。实验结果表明，该路线能够支持多类 Maxwell 任务，并能通过“仿真反馈给大模型再修正”的方式处理部分约束冲突。当前成果适合作为工业软件智能体化项目的阶段性原型，后续可在任务覆盖范围和工程精度上继续扩展。

\end{document}
